%=========================================================
% C++ Chapter 6 Lab Handout (Verbatim Korean-safe)
% Topic: Function Overloading & static Members
% Compile with XeLaTeX
%=========================================================
\documentclass[11pt,a4paper]{article}

\usepackage[margin=1in]{geometry}
\usepackage{fontspec}
\usepackage{kotex}
\usepackage{hyperref}
\usepackage{enumitem}
\usepackage{fancyhdr}
\usepackage{titlesec}
\usepackage{xcolor}
\usepackage{tcolorbox}
\usepackage{fvextra}

% ---------- Fonts ----------
% If D2Coding isn't available, try: NanumGothicCoding, Noto Sans Mono CJK KR
\setmonofont{D2Coding}[Scale=MatchLowercase]

% ---------- Colors ----------
\definecolor{CodeBg}{RGB}{248,248,248}
\definecolor{CodeFrame}{RGB}{220,220,220}
\definecolor{Accent}{RGB}{40,80,160}

% ---------- Header ----------
\pagestyle{fancy}
\fancyhf{}
\lhead{\textbf{C++ 6장 실습}}
\rhead{\textbf{함수 중복과 static 멤버}}
\cfoot{\thepage}

\titleformat{\section}{\Large\bfseries\color{Accent}}{\thesection}{0.6em}{}
\titleformat{\subsection}{\large\bfseries}{\thesubsection}{0.6em}{}

% ---------- Boxes ----------
\newtcolorbox{GoalBox}{
  colback=CodeBg,
  colframe=CodeFrame,
  arc=3pt,
  boxrule=0.8pt,
  left=10pt,right=10pt,top=8pt,bottom=8pt
}
\newcommand{\filetag}[1]{\texttt{\color{Accent}#1}}

% ---------- Code block environment (Unicode-safe) ----------
\DefineVerbatimEnvironment{cppcode}{Verbatim}{
  fontsize=\small,
  frame=single,
  framerule=0.6pt,
  rulecolor=\color{CodeFrame},
  numbers=left,
  numbersep=8pt,
  baselinestretch=1,
  breaklines=true,
  breaksymbolleft=,
  breaksymbolright=,
  commandchars=\\\{\},
  xleftmargin=1.5em
}

\begin{document}

\begin{center}
  {\LARGE \textbf{C++ 6장 실습: 함수 중복과 static 멤버}}\\
  \vspace{6pt}
  {\large (overloading / default parameter / ambiguity / static member variable \& function)}
\end{center}

\vspace{8pt}

\section*{학습 목표}
\begin{GoalBox}
\begin{enumerate}[leftmargin=18pt]
  \item 함수 중복의 개념을 이해하고, 중복 함수를 작성할 수 있다.
  \item 디폴트 매개 변수를 이해하고 작성할 수 있다.
  \item 함수 중복 시 발생하는 모호성의 경우를 판별할 수 있다.
  \item static 속성으로 선언된 멤버의 특성을 이해하고, static 속성을 활용할 수 있다.
\end{enumerate}
\end{GoalBox}

\section*{실습 운영 규칙}
\begin{itemize}[leftmargin=18pt]
  \item 각 실습은 보통 \texttt{main()}을 포함하여 \textbf{독립 실행} 가능합니다.
  \item \textbf{오버로딩 실패/모호성} 예제는 일부러 컴파일 에러가 나도록 만든 코드가 있습니다.
        이런 코드는 주석으로 \texttt{// (에러 확인)} 표기했습니다.
  \item static 멤버 변수는 \textbf{클래스 밖에서 1회 정의}가 필요합니다(링크 오류 방지).
\end{itemize}

\tableofcontents
\newpage

\section{함수 중복 기본: ``이름 같고 매개변수 다르면 OK''}
\begin{GoalBox}
\textbf{핵심}
\begin{itemize}[leftmargin=18pt]
  \item 함수 중복(오버로딩): \textbf{같은 이름}의 함수가 공존.
  \item 구분 기준: \textbf{매개변수의 타입/개수}. \textbf{리턴 타입은 중복과 무관}.
\end{itemize}
\end{GoalBox}

\noindent\filetag{L1\_SumOverloadOK.cpp}
\begin{cppcode}
#include <iostream>
using namespace std;

int sum(int a, int b) { return a + b; }
int sum(int a, int b, int c) { return a + b + c; }
double sum(double a, double b) { return a + b; }

int main() {
    cout << sum(2, 6) << "\n";
    cout << sum(2, 5, 33) << "\n";
    cout << sum(12.5, 33.6) << "\n";
}
\end{cppcode}

\section{오버로딩 실패 1: 리턴 타입만 다르면 안 됨}
\begin{GoalBox}
\textbf{핵심}
\begin{itemize}[leftmargin=18pt]
  \item 아래 두 함수는 매개변수가 완전히 같아서 \textbf{호출 시 구분 불가}.
  \item 따라서 ``리턴 타입만 다르게'' 해서 오버로딩은 \textbf{불가능}.
\end{itemize}
\end{GoalBox}

\noindent\filetag{L2\_OverloadFailReturnType.cpp}
\begin{cppcode}
#include <iostream>
using namespace std;

// (에러 확인) 리턴 타입만 다르면 오버로딩 불가
int sum(int a, int b) { return a + b; }
double sum(int a, int b) { return (double)(a + b); }

int main() {
    cout << sum(2, 5) << "\n";
}
\end{cppcode}

\section{예제 6-1: big() 오버로딩(값 vs 배열)}
\begin{GoalBox}
\textbf{핵심}
\begin{itemize}[leftmargin=18pt]
  \item 같은 이름(big)으로 ``두 정수 중 큰 값''과 ``배열에서 최대값''을 제공.
  \item 호출 형태가 달라서 컴파일러가 자동 구분.
\end{itemize}
\end{GoalBox}

\noindent\filetag{L3\_BigOverload.cpp}
\begin{cppcode}
#include <iostream>
using namespace std;

int big(int a, int b) {
    return (a > b) ? a : b;
}

int big(int a[], int size) {
    int res = a[0];
    for (int i = 1; i < size; i++)
        if (res < a[i]) res = a[i];
    return res;
}

int main() {
    int array[5] = {1, 9, -2, 8, 6};
    cout << big(2, 3) << "\n";
    cout << big(array, 5) << "\n";
}
\end{cppcode}

\section{예제 6-2: sum() 오버로딩(구간합 vs 0부터)}
\begin{GoalBox}
\textbf{핵심}
\begin{itemize}[leftmargin=18pt]
  \item \texttt{sum(a,b)}: a~b 합
  \item \texttt{sum(a)}: 0~a 합
\end{itemize}
\end{GoalBox}

\noindent\filetag{L4\_SumRangeOverload.cpp}
\begin{cppcode}
#include <iostream>
using namespace std;

int sum(int a, int b) {
    int s = 0;
    for (int i = a; i <= b; i++) s += i;
    return s;
}

int sum(int a) {
    int s = 0;
    for (int i = 0; i <= a; i++) s += i;
    return s;
}

int main() {
    cout << sum(3, 5) << "\n";
    cout << sum(3) << "\n";
    cout << sum(100) << "\n";
}
\end{cppcode}

\section{디폴트 매개 변수: ``값이 안 오면 기본값''}
\begin{GoalBox}
\textbf{핵심}
\begin{itemize}[leftmargin=18pt]
  \item 선언 형태: \texttt{type param = defaultValue}
  \item 호출에서 해당 인수를 생략하면 기본값이 자동 전달됨.
\end{itemize}
\end{GoalBox}

\noindent\filetag{L5\_DefaultParamBasic.cpp}
\begin{cppcode}
#include <iostream>
#include <string>
using namespace std;

void star(int a = 5);
void msg(int id, string text = "");

void star(int a) {
    for (int i = 0; i < a; i++) cout << '*';
    cout << "\n";
}
void msg(int id, string text) {
    cout << id << " " << text << "\n";
}

int main() {
    star();
    star(10);
    msg(10);
    msg(10, "Hello");
}
\end{cppcode}

\section{디폴트 매개 변수 제약: ``뒤쪽부터''만 가능}
\begin{GoalBox}
\textbf{핵심}
\begin{itemize}[leftmargin=18pt]
  \item 디폴트 매개 변수는 \textbf{끝 쪽에 몰려} 선언되어야 함.
\end{itemize}
\end{GoalBox}

\noindent\filetag{L6\_DefaultParamRules.cpp}
\begin{cppcode}
#include <iostream>
using namespace std;

// void calc(int a, int b=5, int c, int d=0); // (에러 확인)

void calc(int a, int b = 5, int c = 0, int d = 0) {
    cout << a << "," << b << "," << c << "," << d << "\n";
}

int main() {
    calc(10);
    calc(10, 5);
    calc(10, 5, 20);
    calc(10, 5, 20, 30);
}
\end{cppcode}

\section{오버로딩/디폴트로 인한 모호성}
\begin{GoalBox}
\textbf{핵심}
\begin{itemize}[leftmargin=18pt]
  \item 중복 함수와 디폴트 매개변수 함수를 함께 두면 호출이 모호해질 수 있음.
\end{itemize}
\end{GoalBox}

\noindent\filetag{L7\_AmbiguityDefaultParam.cpp}
\begin{cppcode}
#include <iostream>
#include <string>
using namespace std;

void msg(int id) { cout << id << "\n"; }
void msg(int id, string s = "") { cout << id << ":" << s << "\n"; }

int main() {
    msg(5, "Good Morning");
    // msg(6); // (에러 확인) 호출 모호
}
\end{cppcode}

\section{형 변환으로 인한 모호성}
\begin{GoalBox}
\textbf{핵심}
\begin{itemize}[leftmargin=18pt]
  \item int 인수는 float/double 중 어디로 변환할지 모호할 수 있음.
\end{itemize}
\end{GoalBox}

\noindent\filetag{L8\_AmbiguityCast.cpp}
\begin{cppcode}
#include <iostream>
using namespace std;

float square(float a) { return a * a; }
double square(double a) { return a * a; }

int main() {
    cout << square(3.0) << "\n";
    // cout << square(3) << "\n"; // (에러 확인)
}
\end{cppcode}

\section{참조 매개변수로 인한 모호성}
\begin{GoalBox}
\textbf{핵심}
\begin{itemize}[leftmargin=18pt]
  \item 값/참조 오버로딩은 호출이 구분되지 않아 모호해질 수 있음.
\end{itemize}
\end{GoalBox}

\noindent\filetag{L9\_AmbiguityRefParam.cpp}
\begin{cppcode}
#include <iostream>
using namespace std;

int add(int a, int b) { return a + b; }
int add(int a, int& b) { b = b + a; return b; }

int main() {
    int s = 10, t = 20;
    // cout << add(s, t) << "\n"; // (에러 확인)
}
\end{cppcode}

\section{디폴트로 중복 간소화: fillLine()}
\begin{GoalBox}
\textbf{핵심}
\begin{itemize}[leftmargin=18pt]
  \item 디폴트 매개변수로 여러 중복 함수를 1개로 합칠 수 있음.
\end{itemize}
\end{GoalBox}

\noindent\filetag{L10\_FillLineDefault.cpp}
\begin{cppcode}
#include <iostream>
using namespace std;

void fillLine(int n = 25, char c = '*') {
    for (int i = 0; i < n; i++) cout << c;
    cout << "\n";
}

int main() {
    fillLine();
    fillLine(10, '%');
}
\end{cppcode}

\section{생성자 중복 간소화: MyVector(int n=100)}
\begin{GoalBox}
\textbf{핵심}
\begin{itemize}[leftmargin=18pt]
  \item 기본 생성자/매개변수 생성자를 디폴트 매개변수 생성자 1개로 간소화.
\end{itemize}
\end{GoalBox}

\noindent\filetag{L11\_MyVectorCtorDefault.cpp}
\begin{cppcode}
#include <iostream>
using namespace std;

class MyVector {
    int* p;
    int size;
public:
    MyVector(int n = 100) {
        p = new int[n];
        size = n;
    }
    ~MyVector() { delete[] p; }
    int getSize() const { return size; }
};

int main() {
    MyVector* v1 = new MyVector();
    MyVector* v2 = new MyVector(1024);
    cout << v1->getSize() << "\n";
    cout << v2->getSize() << "\n";
    delete v1;
    delete v2;
}
\end{cppcode}

\section{static 멤버 기본 + 외부 정의(링크 오류 방지)}
\begin{GoalBox}
\textbf{핵심}
\begin{itemize}[leftmargin=18pt]
  \item static 멤버 변수는 클래스 밖에서 1회 정의해야 함.
\end{itemize}
\end{GoalBox}

\noindent\filetag{L12\_StaticMemberBasic.cpp}
\begin{cppcode}
#include <iostream>
using namespace std;

class Person {
public:
    int money;
    void addMoney(int money) { this->money += money; }

    static int sharedMoney;
    static void addShared(int n) { sharedMoney += n; }
};

int Person::sharedMoney = 10;

int main() {
    Person han;
    han.money = 100;
    han.sharedMoney = 200;

    Person lee;
    lee.money = 150;
    lee.addMoney(200);
    lee.addShared(200);

    cout << han.money << " " << lee.money << "\n";
    cout << han.sharedMoney << " " << lee.sharedMoney << "\n";
}
\end{cppcode}

\section{static 접근: 클래스명:: vs 객체}
\begin{GoalBox}
\textbf{핵심}
\begin{itemize}[leftmargin=18pt]
  \item static은 \texttt{ClassName::}로 접근 가능, non-static은 불가.
\end{itemize}
\end{GoalBox}

\noindent\filetag{L13\_StaticAccessClassScope.cpp}
\begin{cppcode}
#include <iostream>
using namespace std;

class Person {
public:
    double money;
    static int sharedMoney;
    static void addShared(int n) { sharedMoney += n; }
};
int Person::sharedMoney = 10;

int main() {
    Person::addShared(50);
    cout << Person::sharedMoney << "\n";

    Person han;
    han.money = 100;
    han.sharedMoney = 200;

    Person::sharedMoney = 300;
    Person::addShared(100);

    cout << han.money << " " << Person::sharedMoney << "\n";
}
\end{cppcode}

\section{static으로 전역 함수 캡슐화: Math}
\begin{GoalBox}
\textbf{핵심}
\begin{itemize}[leftmargin=18pt]
  \item 관련 전역 함수를 클래스로 묶어 사용성을 개선.
\end{itemize}
\end{GoalBox}

\noindent\filetag{L14\_MathStaticClass.cpp}
\begin{cppcode}
#include <iostream>
using namespace std;

class Math {
public:
    static int abs(int a) { return (a > 0) ? a : -a; }
    static int max(int a, int b) { return (a > b) ? a : b; }
    static int min(int a, int b) { return (a > b) ? b : a; }
};

int main() {
    cout << Math::abs(-5) << "\n";
    cout << Math::max(10, 8) << "\n";
    cout << Math::min(-3, -8) << "\n";
}
\end{cppcode}

\section{static으로 객체 개수 세기: numOfCircles}
\begin{GoalBox}
\textbf{핵심}
\begin{itemize}[leftmargin=18pt]
  \item new[] / delete[]가 생성자/소멸자 호출 횟수에 영향을 준다.
\end{itemize}
\end{GoalBox}

\noindent\filetag{L15\_CountObjectsStatic.cpp}
\begin{cppcode}
#include <iostream>
using namespace std;

class Circle {
private:
    static int numOfCircles;
    int radius;
public:
    Circle(int r = 1) : radius(r) { numOfCircles++; }
    ~Circle() { numOfCircles--; }
    static int getNumOfCircles() { return numOfCircles; }
};
int Circle::numOfCircles = 0;

int main() {
    Circle* p = new Circle[10];
    cout << "생존하고 있는 원의 개수 = " << Circle::getNumOfCircles() << "\n";
    delete[] p;
    cout << "생존하고 있는 원의 개수 = " << Circle::getNumOfCircles() << "\n";

    Circle a;
    cout << "생존하고 있는 원의 개수 = " << Circle::getNumOfCircles() << "\n";
    Circle b;
    cout << "생존하고 있는 원의 개수 = " << Circle::getNumOfCircles() << "\n";
}
\end{cppcode}

\section{static 함수 제한: this/non-static 접근 불가}
\begin{GoalBox}
\textbf{핵심}
\begin{itemize}[leftmargin=18pt]
  \item static 함수는 객체가 없어도 호출 가능 $\rightarrow$ this 없음.
\end{itemize}
\end{GoalBox}

\noindent\filetag{L16\_StaticCannotAccessThis.cpp}
\begin{cppcode}
#include <iostream>
using namespace std;

class PersonError {
    int money;
public:
    // static int getMoney() { return money; } // (에러 확인)
    void setMoney(int money) { this->money = money; }
};

int main() {
    PersonError e;
    e.setMoney(100);
    cout << "OK\n";
}
\end{cppcode}

\section*{마무리 체크리스트}
\begin{GoalBox}
\begin{itemize}[leftmargin=18pt]
  \item 오버로딩 성공/실패 조건을 정확히 말할 수 있는가?
  \item 디폴트 매개변수 규칙(뒤쪽부터)과 모호성 사례를 설명할 수 있는가?
  \item static 멤버 변수 외부 정의가 왜 필요한가(링크 오류)?
  \item static 멤버 함수의 제한(this/non-static 접근 불가)을 설명할 수 있는가?
\end{itemize}
\end{GoalBox}

\end{document}
